%!TEX program = lualatex
\documentclass[ngerman,headheight=70pt]{scrartcl}
\usepackage[ngerman]{babel}
\usepackage{blindtext}
\usepackage[utf8]{luainputenc}
\usepackage{scrlayer-scrpage}
\usepackage[T1]{fontenc}
\usepackage{fontspec}
\usepackage{multicol}
\usepackage{enumitem}
\usepackage[german=quotes]{csquotes}
\usepackage[bottom=2cm,footskip=8mm]{geometry}
\usepackage{pdfpages}
\setmainfont{Source Sans Pro}
\MakeOuterQuote{"}

\setlength{\parindent}{0pt}
\setlength{\parskip}{6pt}

\pagestyle{plain.scrheadings}
\thispagestyle{scrheadings}
\lohead{\vspace*{-1.5cm}\includegraphics[scale=0.4]{../images/stupa_logo_rot.pdf}}
\rohead{\\Präsidium des\\Studierendenparlaments}

\lofoot[Universität Hamburg · Präsidium des Studierendenparlaments\\\hrule Von-Melle-Park 5 · D-20146 Hamburg]{Universität Hamburg · Präsidium des Studierendenparlaments\\\hrule Von-Melle-Park 5 · D-20146 Hamburg}
\cofoot[]{}
\cohead[\thepage]{}

\renewcommand{\thesubsection}{TOP \arabic{subsection}}
\renewcommand{\thesubsubsection}{\arabic{subsubsection}.}
\setcounter{subsection}{-1}

\begin{document}
    UHH · StuPa-Präsidium · Von-Melle-Park 5 · D-20146 Hamburg

    %%% PAGE ONE %%%%
    \section*{Protokoll der 1. Fortsetzung der 1. Sitzung des Studierendenparlaments vom 28. April 2016}

    \textbf{Protokoll: Jim Martens}\\
    \textbf{Ort: Raum S08 in der ehemaligen HWP }\\
    \textbf{Beginn: 18.25 Uhr}\\
    \textbf{Ende: 23.41 Uhr}

    \vspace{0.5cm}
    \begin{tabular}{ll}
        Anwesend: & \\
            RCDS (5 Sitze): & Ramon Weilinger, Antonia Niecke, Ramin Shakiba, \\
                            & Jennifer Maack, Benjamin Welling \\
             CampusGrün (14 Sitze): & Laura Franzen, Geoffrey Youett, Elena Rysikova, \\
                                   & Philipp Droll, Yasemin Günther (später), Melf Johannsen,\\
                                   & Tahnee Herzig, Mario Moldenhauer, Freya Schmitz, \\
                                   & Jim Martens, Svenja Horn, Mirzo Ulugbek Khatamov \\
                                   & Martin Sievert, Armin Günther \\
             Bier-Liste (2 Sitze) : & Karen Martirosian (Bier) \\
             WiWi (2 Sitze): & Claas-Friso Hente, Kay Zöllmer \\
             Unicorns (5 Sitze): & Katharina Kucza, Johannes Peplow, Annkathrin Löffler, \\
                                 & Andreas Hartkamp \\
             Liste LINKS (3 Sitze): & Gunhild Berdal, Till Petersen, Sinah Mielich \\
             HWP (2 Sitze): & Samet Gunay, Ajdina Karahasan \\
             MIN (4 Sitze): &  \\
             SDS* (3 Sitze): & Mena Winkler, Jacob Petersein, Artur Brückmann \\
             Bart-LISTE (2 Sitze): & Dominic Laumer \\
             LHG (1 Sitz): & Tobias Heisig \\
             harte zeiten (1 Sitz): & Tobias Berking \\
             Jura (1 Sitz): &  \\
             AL (2 Sitze): & Karima Schulze, Henri Weber \\
            & \\
        Entschuldigt: & Lotte Rullkötter, Timo Zeimet, Thea Wahlers \\
                      & Jan Detampel, Jakob Pape, Ailina Salten\\
        Unentschuldigt abwesend: &  Marielle Hermstrüwer, Johann Baumhoefener (Jura)\\
                                &\\
        Rücktritte: & \\
    \end{tabular}
    \newpage
    %%% Vorgeschlagene Tagesordnung %%%
    \underline{Vorgeschlagene Tagesordnung}
    \begin{enumerate}[label={\textbf{Top \theenumi}},leftmargin=*]
        \setcounter{enumi}{10}
        \item Bestätigung des RBCS-Referenten (V1617-010)
        \item \#uhh hilft (V1617-007,V1617-012)
        \item
            \begin{enumerate}
                \item Rechenschaftsbericht des amtierenden AStA
                \item Fragen und Diskussion
                \item Entlastung des AStA
            \end{enumerate}
        \item Wahl des neuen AStA-Vorstandes
            \begin{enumerate}
                \item Diskussion VS-Thesen (V1617-011)
                \item Wahl des AStA-Vorstandes
            \end{enumerate}
        \item Bestätigung der AStA-Referent*innen
        \item Antrag Ältestenrat zur Beauftragung des SWOGA (V1617-013)
        \item Verschiedenes
    \end{enumerate}

    \newpage

    {\Large \underline{Teil A}}

    \subsection{Formalia}

    \subsubsection{Geschäftsbericht des Präsidiums}

    Ramon berichtet, dass das Präsidium sich als gleichberechtigt wahrnimmt und
    daher die Posten Präsident*In, Vizepräsident*In und Schriftführer*In nicht
    verteilt hat.

    Ferner erklärt Ramon, dass das Präsidium in Bezug auf den Verlust des Mandates
    einen Sitzungstag und nicht etwa eine Sitzung als das wichtige
    Entscheidungskriterium ansieht.

    Ausschusskonstituierungen sind durch und Ergebnisse können auf der StuPa-Website
    eingesehen werden. Sofern die Zusammenarbeit mit den Ausschüssen funktioniert,
    werden auch in Zukunft die Protokolle auf der Website erscheinen.

    \subsubsection{Feststellung der endgültigen Fassung des Teils B der Tagesordnung}

    Ramon schlägt vor, dass die Bestätigung des RBCS-Referenten dann stattfindet,
    wenn dieser anwesend ist und der dann stattfindene TOP fertig ist.

    Till fragt nach, ob TOP 13 als einheitlicher TOP gesehen wird. Ramon antwortet,
    dass das Präsidium dies so sieht.

    Danach kein Widerspruch zum Vorschlag.

    %%% BESCHLOSSENE TAGESORDNUNG %%%
    \underline{Endgültige Fassung der Tagesordnung}
    \begin{enumerate}[label={\textbf{Top \theenumi}},leftmargin=*]
        \setcounter{enumi}{10}
        \item Bestätigung des RBCS-Referenten (V1617-010)
        \item \#uhh hilft (V1617-007,V1617-012)
        \item
            \begin{enumerate}
                \item Rechenschaftsbericht des amtierenden AStA
                \item Fragen und Diskussion
                \item Entlastung des AStA
            \end{enumerate}
        \item Wahl des neuen AStA-Vorstandes
            \begin{enumerate}
                \item Diskussion VS-Thesen (V1617-011)
                \item Wahl des AStA-Vorstandes
            \end{enumerate}
        \item Bestätigung der AStA-Referent*innen
        \item Antrag Ältestenrat zur Beauftragung des SWOGA (V1617-013)
        \item Verschiedenes
    \end{enumerate}

    \setcounter{subsection}{10}
    \vspace{1cm}
    {\Large \underline{Teil B}}

    \subsection{Bestätigung des RBCS-Referenten}

    \subsubsection{Bestätigung der Wahlniederschrift}

    Geoffrey stellt die Wahlniederschrift vor. Es gab leider eine geringe
    Wahlbeteiligung.

    Keine Fragen zum Protokoll.

    M:0:0, Protokoll angenommen.

    \subsubsection{Bestätigung des Referenten}

    Lars stellt sich vor. Wahlbeteiligung ist seit längerer Zeit gering.
    Problem ist, dass die Gruppe sehr divers und schwer zu zentralisieren
    ist.

    Beratung wird stark frequentiert. Lars hat im vergangenen Semester
    einen Schwerpunkt auf Beratung gelegt. Bisher zwei Beraterinnen,
    die auf eigenen Wunsch ehrenamtlich arbeiten. RBCS versucht mehr
    Projektkräfte zu bekommen, um mehr Veranstaltungen anbieten zu können.

    Ergonomiekurs war ausgebucht. Selbstverteidigung für Frauen war jedoch
    von nur 3 Teilnehmerinnen besucht.

    Dieses Semester ist die größte Veranstaltung der Berater*Innen-Workshop
    für ungefähr 10 Studierende. Workshop findet an einem Wochenende
    Anfang Juni statt. Nächsten Monat findet auch Ergonomieseminar noch einmal
    statt.

    Till sagt, dass Beratung und "Stellvertreterpolitik" Grund für geringe
    Beteiligung sind. Hat das RBCS vor in Zukunft mehr politische Arbeit
    zu machen?

    Artur sagt, dass aktuelle politisch-kulturelle Situation dazu führt, dass
    sich Leute ungern der Gruppe BCS zuordnen.

    Golnar fragt, ob es sinnvoll wäre sich mit der Geschichte in dem Bereich
    BCS auseinanderzusetzen.

    Lars sagt, dass er es ähnlich sieht, wie Till, Artur und Golnar. Das RBCS
    setzt sich durchaus auch politisch ein. Mehr politische Arbeit
    geplant, aber immer schwer Leute zu erreichen.

    21:00

    M:0:0 für Bestätigung

    \subsection{\#uhh hilft}

    Golnar stellt den Antrag V1617-007 zu \#uhh hilft vor.

    Jennifer stellt den Ersetzungsantrag V1617-012 zu \#uhh hilft vor.

    Geoffrey fragt aus welcher Motivation heraus der Ersetzungsantrag kommt.
    Ihm fehlt die Situationsanalyse in dem Ersetzungsantrag und fragt die
    Antragssteller*Innen, ob dies beabsichtigt ist und bittet sie generell
    darauf einzugehen.

    Till erläutert, warum die Situationsanalyse und Begründung wichtig sind und
    ruft dabei Punkte aus der Situationsanalyse und Begründung auf.

    Jennifer erklärt, dass sie das Ergebnis des Antrags gut finden, aber nicht
    die Annahmen und deswegen haben sie einen Ersetzungsantrag gemacht.

    Golnar sagt, dass sich dann die Frage stellt, warum das Ergebnis gut gefunden
    wird, wenn die Annahmen nicht geteilt werden. Sie erläutert, was die inhaltliche
    Ausrichtung der AG Flüchtlinge? ist. Sie sagt, dass eine politische Begründung
    des Antrags notwendig ist, um das Programm auch durchzusetzen.

    Gunhild findet es bemerkenswert, dass die Uni gegen das Prinzip des teile
    und herrsche vorgeht. Sie findet es verlogen zu sagen, dass man die Ziele
    des Antrags unterstützt, aber nicht bereit ist die Politik der Bundesregierung
    inhaltlich zu kritisieren. Sie sagt, dass DaF eine historische Bedeutung
    an der Uni hat und wieder in der Germanistik angesiedelt werden sollte.
    Das Anfängerniveau sollte massiv ausgebaut werden. Sie findet, dass dieser
    Ausbau von DaF allen zugute kommt - gerade vor dem Hintergrund, dass häufig
    Flüchtlinge als Quell vielen Übels gesehen werden.

    Franziska findet es bemerkenswert, dass das Programm die Flüchtlinge in alle
    Facetten der Universität aufnimmt. Sie hat das Programm so verstanden, dass
    ein Ziel ist den Flüchtlingen einen Überblick über die politische Situation
    in Deutschland zu verschaffen. Sie sagt, dass der Ausblick, dass dieses
    Programm nur den Anfang einer Bewegung darstellen kann, in dem Ersetzungsantrag
    fehlt. Man hat eine Position erreicht, dass dieses Programm richtig und gut ist,
    aber die Welt ist noch nicht gut.

    Golnar hat Änderungsantrag eingereicht (siehe Sitzungsordner).

    6:M:1 für Ersetzungsantrag.

    Golnars Änderungsantrag wird von den Antragsstellern übernommen.

    > in Zeile 37 einfügen
    > Darüber hinaus fordert das Stupa die Bürgerschaft auf, den Hamburger
    Hochschulen für den bedarfsgerechten Ausbau von DaF und die Durchführung
    der Studien(vorbereitungs)programme für Geflüchtete
    dauerhaft bedarfsgerechte Mittel zusätzlich zuzuweisen.

    M:0:2 für Gesamtantrag.

    18:57

    Ramon stellt fest, dass Lars noch nicht anwesend ist.

    \subsection{Entlastung des AStA}

    \subsubsection{Rechenschaftsbericht des AStA}

    Moritz hält seinen Rechenschaftsbericht. Er könnte von 2 Legislaturen enthalten.
    Moritz erläutert, wie die Rechenschaftsberichte aufgeteilt sind. Seiner
    wird hauptsächlich aus Vorstandsperspektive stattfinden.

    Vorstand:
    1. Formale Aufgaben
    2. inhaltliche Richtung andiskutieren und vorbereiten (pol. Vorstand)
    3. Repräsentation nach Außen in Uni
    4. Repräsentation nach außen in die Stadt

    1.
    Vorstand verlässt sich auf Festangestellte. Vorstand muss eigentlich nur
    Sonderfälle behandeln. Regelfälle werden "quasi wie von selbst" durch die
    Festangestellten übernommen.

    Zielsetzung die Urlaube so zu legen, dass immer ein Zeichnungsberechtigter
    anwesend war.

    Hoher Zeitaufwand für Rechtsstreitigkeiten.
    Überlegungen:

    1. Eine der Festangestellten (z.B. Maike) sollte sich mit
    Arbeitsrecht beschäftigen. Ansonsten schwer im AStA an Kompetenz für Arbeitsrecht
    zu kommen.
    2. AStA sollte öfter darüber diskutieren, warum die VS etwas tun sollte.
    Daraus ist es häufig einfacher zu sagen, warum die VS etwas tun darf.
    3. Erfahrungen weitergeben.

    2.
    Vorstand hätte eher Konflikte im Vornherein besprechen müssen. Pol. Vorstand
    sollte in Zukunft stärker genutzt werden.

    3.
    Repräsentation war einfacher, da man sich nicht mehr vorstellen musste.
    AStA ist wichtiger Gesprächspartner in der Uni. Kontakt zu den Lehrenden
    sollte weiter ausgebaut werden. Anerkennung wurde erworben und man schätzt
    den AStA. Selbst Leute, die inhaltlich nicht dem AStA zustimmen, sich über
    eine aktive VS freuen.

    4.
    Mit NOlympia einen großen Eindruck hinterlassen. Es waren Schulklassen auf
    dem Campus und der AStA war auch in Schulveranstaltungen. Eigene inhaltliche
    Veranstaltungen sind ausbaufähig. Positives Beispiel: Kapitalismus und Nachhaltigkeit.

    Bau:
    Dem AStA war wichtig die Nutzerinteressen in dem Verfahren zu verankern.
    Bündelung der Interessen und Vermittlung ist ebenfalls wichtig.
    Moritz hat dieser Part am meisten Spaß gemacht.

    -----------------

    Vincent fährt fort mit seinem Rechenschaftsbericht.
    Der Vorstand vertritt den AStA nach innen und außen. Er hat formale und
    politische Verantwortung. Vincent sieht den Finanzvorstand als Kristallationspunkt
    zwischen formalen Vorgaben und inhaltlicher Ausrichtung. Vincent hat viel
    zu Arbeitsrecht und Finanzrichtlinien gelernt.

    Linker AStA eckt auch an, aber das ist wichtig. Bürgerschaftsanfragen und Klagen sind
    nicht immer spaßig, aber teilweise notwendig. Wichtig relevante Rechtsangelegenheiten
    von unwichtigen zu unterscheiden.

    Vincent listet einige Punkte der inhaltlichen Arbeit des AStA auf. Der Vorstand
    hat in vielen Bündnissen mitgearbeitet. Eine engere Zusammenarbeit mit dem fzs
    sollte angestrebt werden.

    Amt als Vorsitzender ist sehr vereinnahmend. Vincent ist dem Anspruch
    zwischen formalen Vorgaben und egalitärer Arbeitsteilung nicht immer
    gerecht geworden.

    ---------------

    pol. Vorstand

    Franziska hält ihren Rechenschaftsbericht. Im vergangenen Jahr wurde vieles
    erreicht. Beispielhaft wird der Dies Academicus erwähnt. AStA Arbeit sollte
    immer wie eine Dies Academicus Vorbereitung sein. Universität sollte jeden
    Tag wie Dies Academicus werden.

    ---------------

    Finanzreferat
    Timo hält seinen Rechenschaftsbericht. Er wurde erst im Januar Referent.
    Sein Bericht kann dadurch nur die letzten 3 Monate umfassen. Mit Anspruch
    angetreten "so inhaltlich wie möglich, so formal wie nötig". Anfangs nur
    theoretisches Wissen. Große Anfrage der FDP hat ihm viel Wissen gebracht.
    Er ist sich nicht sicher, ob das Finanzreferat das kommunikativste Referat
    ist. Timo hat jeden mal gesehen. Alle Töpfe sind im Plus, außer der Topf
    des Präsidiums. Er hat nicht das Gefühl, dass Leute auf der Strecke bleiben
    wegen Bürokratie.

    ---------------

    HoPo-Referat
    Philipp

    Aufgrund der langen Zusammenarbeit wurde sich viel Handlungsfähigkeit
    erarbeitet. NOlympia gutes Beispiel. Gute Möglichkeiten der Zusammenarbeit
    mit anderen Studierendenschaften in Hamburg (z.B. LAK).
    Gute Zusammenarbeit aus Hamburg auf fzs übertragen - gerade vor dem Hintergrund,
    dass Uni Hamburg im bundesweiten Vergleich sehr viel Möglichkeiten hat.

    Mena

    19:30

    Vieles wurde bereits gesagt. Sie findet die Ausrede "ich muss noch lernen"
    sehr schlecht. Lernen sollte etwas Positives sein. Ausrede zeugt von
    degeneriertem Verständnis von Lernen. Rahmenprüfungsordnung ist letztes
    Jahr entwickelt worden und befindet sich im Endprozess. Mit der RPO
    ist Verbesserung der Studienbedingungen möglich.

    Sie sieht Verbesserungspotenzial im AStA bzgl. der grundlegenden inhaltichen
    Ausrichtung.

    Meta: Die redet echt wie Karla Kolumna. Ohne Punkt und Komma.

    Gemeinsames Lernen und Lehren sollte in Zukunft weiter angegangen werden.

    ------------

    Soziales
    Artur

    Artur liest einen Artikel von der ZEIT zu Olympia vor und bewertet diesen.
    Ferner geht er noch einmal auf die Arbeit des AStA zu Olympia ein.
    In Zukunft sollten wissenschaftlich-aufklärerischen Argumente stärker
    in die Gesellschaft getragen werden.

    Artur zitiert die AStA-Leitlinien von 2015 und geht inhaltlich auf die
    genannten Punkte ein.

    Schwierigkeiten dabei Verhältnis zwischen Hilfe im Bestehenden und
    strukturelle Verbesserung gut hinzubekommen. Zu häufig haben Auseinandersetzungen
    z.B. mit dem Studierendenwerk im Hinterzimmer stattgefunden.

    AStA muss insgesamt stärker herauskommen und inhaltliche Arbeit stärker
    verallgemeinern.

    ------------

    Ökologie
    Mario

    Mario hat im Oktober als Quereinsteiger im Referat für Ökologie und
    Nachhaltige Entwicklung begonnen. Anfangs musste er sich erst einmal
    in die inhaltliche Arbeit des Referats einarbeiten und dankt allen,
    die ihm dabei geholfen haben.

    Er hat zusammen mit dem restlichen Referat die bereits geplanten Veranstaltungen
    umgesetzt. Neue Veranstaltungen konnten leider nicht im geplanten
    Umfang durchgeführt werden.

    In Bezug auf seinen Anspruch zieht er ein gemischtes Fazit. Mario findet
    die anstehende Zusammenlegung des Referates sinnvoll.


    ------------

    Kultur
    Stefan Seld

    Seld hält seinen Rechenschaftsbericht. Es gibt ein Radio, bei dem
    Themen des AStA diskutiert werden. Freies Senderkombinat wird das Ding
    genannt. Meta: Werbeplug für sein Radio.

    Ramon schlägt vor jetzt die Beschlussfähigkeit festzustellen, anschließend
    eine Präsidiumspause zu machen. Während der Pause sollen die Rechenschaftsberichte
    von Till Mehrpohl und Carina gedruckt werden, sodass das Parlament diese
    zur Kenntnis nehmen kann.

    Till sagt, dass es sinnvoll wäre die Bericht jetzt schon zu drucken, um sie
    dann in der Pause lesen zu können.

    Moritz schlägt vor, dass er die Berichte ausdrucken kann, da er für die
    Beschlussfähigkeit nicht erheblich ist.

    Da einer der Berichte per LibreOffice nicht gelesen werden kann,
    wird technisch bedingt dem Vorschlag Ramons gefolgt.

    \textit{Beschlussfähigkeit mit 35/47 gegeben.}

    Präsidiumspause bis 20:20.

    Es ergibt sich in der Pause, dass Carinas Rechenschaftsbericht
    auch per Word nicht geöffnet werden kann.

    Oliver Vornfeld liest den Rechenschaftsbericht von Till Mehrpohl vor.

    \subsubsection{Fragen und Diskussion}

    Till stellt eine Frage an Moritz und Vincent: An welchen Stellen habt
    ihr den Widerspruch zwischen formaler Hierarchie und egalitärer
    Arbeitsweise gesehen?

    an Moritz: Wie kann die "Kür" der inhaltlichen Arbeit stärker in der
    Arbeit des Vorstandes stattfinden?

    an Moritz: Ziehst du die Konsequenz die FSRK stärker in den AStA einzubinden?

    an Vincent: Wieso siehst du einen Widerspruch zwischen "Reformismus"
    und "Aktionismus"?

    an Timo: Hast du eine Vorstellung, wie aus dem Umstand des Überblicks
    über alle Referate, etwas sinnvolles gemacht werden kann?

    Moritz beantwortet die ersten beiden Fragen. Die Aufgabe des Vorstandes
    ist es an allen Stellen mitzuarbeiten. Service ist nur solange Sklaverei,
    solange man nicht an inhaltlicher Arbeit beteiligt ist. Strukturelle
    Arbeit wurde im Vorstand gebündelt, weil Arbeitsteilung manchmal auch
    sinnvoll ist.

    Moritz beantwortet die Frage zur FSRK. Beteiligung der FSRe ist in Hamburg
    nicht optimal gelöst. Er macht dies aber nicht an der FSRK fest. Ideale
    Lösung noch unklar, aber klarer Handlungsbedarf.

    20:30

    Vincent hat viel zum Verhältnis gelernt, als er mit Vorstand angefangen hat.
    Gerade das Verhältnis zu den Festangestellten ist dabei zu nennen.
    Außerdem ist zwischen Amt und Person zu unterscheiden. Der Vorstand ist als
    Amt Arbeitgeber der Festangestellten, nicht die Personen im Vorstand.

    Vincent beantwortet die Frage zum Widerspruch.

    Timo beantwortet die Frage zur Übersicht, die man als Finanzreferent*In
    hat.

    Gunhild fragt, wo die Gewichtung liegt (Anerkennung oder Anecken).

    Till sagt, dass das Finanzreferat der Ort ist, wo potentielle Konflikte
    oder Herausforderungen rechtzeitig thematisiert werden kann. (Reword)
    Bei allem, was man tut, geht es um Menschen. Gremien sind daher kein
    Selbstzweck, sondern für Menschen da.

    Moritz kann die Frage von Gunhild nicht beantworten, da er den Widerspruch
    nicht sieht. Anerkennung erarbeiten durch sinnvolles Maß an Anecken.
    Bei Diskussion um Ordnung macht es keinen Sinn pauschal gegen eine Ordnung
    zu sein.

    Vincent meint, dass es sinnvoll ist immer eine Gesprächsgrundlage mit
    Leuten zu haben.

    20:46

    \subsubsection{Entlastung des AStA}

    M:e:3, AStA damit entlastet

    \subsection{Wahl des neuen AStA-Vorstandes}

    \subsubsection{Diskussion VS-Thesen}

    Franziska macht einen Aufschlag zur Diskussion und stellt
    den Antrag vor.

    Philipp geht auf die Bemühungen um Kapazitätsrecht und Exzellenzinitiative ein.
    Er geht auf die historische Entwicklung um Wissenschaftsauffassung ein.

    Es soll keine einzelnen Vorstellungen geben.

    Fraktionspause RCDS 10 Minuten bis 21:25.
    Dynamisch genommene Präsidiumspause bis 21:35

    Mena, Geoffrey und Janis gehen nach vorne. Mena sagt, dass wir uns in einer
    Entscheidungssituation befinden. Diese Situation soll sich auch auf Aktivitäten
    im Hochschulumfeld auswirken. Dies ist im VS-Thesentext so beschrieben.

    HoPo, Kultur und Öffentlichkeitsarbeit werden thematisch zusammengelegt.
    Universität als Lebensraum. Universität ist ein Lebensmittelpunkt, in dem
    man sich begegnet, sich verändert und letztendlich auch die Welt verändert.

    Janis sagt etwas zum kulturpolitischen Bereich. Wichtig ist, dass sich
    der AStA nicht als Dienstleister sondern als Impulsgeber versteht.
    Der AStA kann insofern als Katalysator gesehen werden und die zentrale
    Stelle nutzen, um alle Bestrebungen zu fördern.

    Geoffrey stellt die Öffentlichkeitsarbeit vor. Ein Paradebeispiel dafür
    ist die Holzwirtschaft, die davor steht abgewickelt zu werden. Es muss
    daraufhin gewirkt werden, dass das Silodenken ("mein Studium, meine Fakultät")
    überwunden wird. In dem Sinne erinnert Geoffrey an die Demo am 3.5., wo es
    nicht nur um die Holzwirtschaft, sondern um die Finanzlage an der Uni insgesamt
    geht. Eine gemeinsame Position entsteht nicht durch eine Website, sondern
    durch das gemeinsame Gespräch.

    Seld hebt hervor, dass es in dem Thesenpapier um Solidarität geht.
    Er liest in dem Papier auch, dass sich die verschiedenen Kämpfe miteinander
    verbinden lassen.

    21:45

    Seld sagt, dass wir in einer Umwelt leben, die stark von Menschen beeinflusst
    wird. Meta: Sein Redebeitrag soll den Cluster Soziales, Ökologie und Soziale
    Bewegungen vorstellen.

    Seld will eine Kampagne anstoßen, um den HVV kostenlos zu machen.

    Franziska hält einen Redebeitrag. Mit dem Text wurde eine wichtige Analyse
    getroffen, die auch für die zukünftige Arbeit relevant ist.

    Jacob Petersein geht auf den "antifaschistischen Impetus" des Textes
    ein.

    Till geht auf den letzten Absatz des Textes ein. Er stellt die Radikalität
    des Vorsatzes dar, die Solidarität auszubauen. Er geht auf zwei Personen
    ein, die früher Teil eines ausgesprochen antilinken AStAs waren, aber erkannt
    haben, dass die Arbeit in der VS extrem befreiend ist. Till findet es wichtig
    zu reflektieren, wie wichtig Anerkennung sein darf, da es den Anerkennenden
    eine Institution zuspricht. Die Alternative zur kapitalistischen Gestaltung
    der Umgebung muss die menschenwürdige Gestaltung der Umgebung sein. Eine
    Absage zu der Dialektik (Anecken vs. Anerkennung) sorgt dafür, dass
    man in eine reaktionäre Argumentation abrutschen kann. Dies sollte vom
    AStA berücksichtigt werden.

    Armin findet, dass in Zukunft eine größere Rolle bei den ökologischen
    Fragen eingenommen werden sollte.

    22:00

    Artur sagt, dass man sich jeglicher rechter Hetze zur Wehr setzen muss.
    Studierende lassen sich zu häufig noch einreden, dass man privilegiert
    sei und das Engagement in AStA und Co nur vorübergehend und eine Phase
    sei. Er geht auf den US-Wahlkampf mit Donald Trump und Bernie Sanders ein.
    Studierende sollten es als gemeinsamen Kampf verstehen, wenn z.B. die Krankenschwestern
    von einem Krankenhaus streiken.
    Artur verweist auch noch einmal auf NOlympia und die laufenden Anstrengungen
    die Arbeiten ins Italienische zu übertragen, damit die Errungenschaften
    in Hamburg auch auf Rom übertragen werden können. Dort soll es jetzt auch
    ein Referendum zu Olympia geben.

    Franziska hält einen weiteren Redebeitrag. Sie geht auf die Beeinflussbarkeit
    der Hamburger Regierung ein und erwähnt die Enthaltung von Fegebank
    bei der Exzellenzinitiative.

    M:e:0, Text angenommen.

    Persönliche Erklärung von Ramin, Jennifer, Claas-Friso, ...
    Wir halten die Diskussion der VS-Thesen wie vorgetragen für die Vorstellung
    des jeweiligen Parteiprogramms und fordern, dass sie auch dementsprechend
    benannt werden.

    Persönliche Erklärung von Geoffrey:
    Die Opposition wieder verkennt
    die Wichtigkeit vom Argument,
    denn will man wirklich überzeugen
    hilfts nicht sich unterm Tisch zu beugen.
    Denn nur durchs Reden, Diskutieren
    wird inhaltlich etwas passieren.

    Eine Erklärung zum Finale
    wird gern auch Klingelstreich genannt.
    Man klopft und lächelt etwas fahle
    doch geht die Tür auf wird gerannt.

    \subsubsection{Wahl des AStA-Vorstandes}

    Es gibt eine Kandidaturliste:

    1. Vorstand: Philipp Droll
    2. Vorstand Franziska Hildebrandt

    Philipp, Franziska und Timo stellen sich vor.

    Till stellt die Frage an Philipp, warum ausgerechnet
    Philipp Vorstand macht.

    Till stellt eine Frage an Philipp und Franziska, welche Punkte sie
    von dem Rechenschaftsbericht von Vincent und Moritz übernommen haben.

    Philipp hofft, dass es noch andere Gründe gibt, als dass er bei CG ist.
    Er hat die Arbeit im AStA sehr positiv wahrgenommen. Vor dem Hintergrund
    der Erkenntnisse von Philipp, sieht er als sinnvoll an, dass er Vorstand
    macht. Er folgt tendenziell der Empfehlung von Vincent, stärker im fzs
    zu arbeiten.

    Franziska geht in vielen Punkten auf die Arbeit des AStA ein.

    Es gibt keine Personaldebatte.

    Die Wahl findet offen statt.

    28:11:0, damit AStA-Vorstand gewählt, Wahl wird angenommen.

    \subsection{Bestätigung der AStA-Referent*Innen}

    Finanzreferent:
    Timo Tomanek

    HoPo, Öffentlichkeitsarbeit, Kultur:
    Mena Winkler, Geoffrey Youett, Janis Franken, Kim Wiesenberg

    Soziales, Ökologie, Soziale Bewegung:
    Artur Brückmann, Armin Günther, Stefan Seld

    Antifaschismus, Antidiskriminierung, Kritische Wissenschaft:
    Karima Schulze, Mario Moldenhauer, Annkathrin Löffler

    22:30

    GO-Antrag von Melf auf Verlängerung bis nach TOP 17. Kein Widerspruch.

    Vorstellung der Referent*Innen. Kim Wiesenberg ist aus gesundheitlichen
    Gründen leider nicht anwesend. Janis stellt Kim vor.

    Till stellt die Frage, an welcher Stelle in der Referatsstruktur die
    Beschäftigung mit der Friedensentwicklung verortet sein soll.

    Till stellt die Frage, welche Bedeutung die Personen aus Antifa und Antidis
    dem sozialen Konflikt zumessen.

    Till findet, dass die Wörter "halt" und "irgendwie" dafür sprechen, dass sich
    Seld cool findet. Gibt es Punkte, die Seld wirklich inhaltlich teilt?

    Golnar hat die Frage, was die Referent*Innen unbedingt anpacken wollen.

    Seld hat Probleme mit dem Thesenpapier, v.a. mit dem Duktus und der
    Entstehungsgeschichte.
    Meta: Er wäscht mal wieder schlechte Wäsche und spricht erneut die
    Schuldenbremse an. Starker Applaus.

    Annkathrin geht auf Tills Frage zu Antifa und Antidis ein. Sie sieht
    den sozialen Konflikt auch als Problem an.

    23:00

    Artur meint, dass Frieden ein Querschnittsthema sein sollte. AStA sollte
    den imperialistischen Impuls zurückweisen. Artur will unbedingt ändern,
    dass Studierenden eingeredet wird, dass sie an den neoliberalen Zuständen
    Schuld seien.

    Mena sagt, dass die Friedensfrage überall aufgehoben sein sollte.
    Sie möchte daran arbeiten, dass weniger nebeneinander her studiert und
    gearbeitet wird.

    Geoffrey sagt auch, dass Frieden übergreifend bearbeitet werden muss.
    Strukturell soll die Diskussion zur Zivilklausel weiter vorangetrieben werden.
    Kein Mensch sollte unter Kriegsbedingungen leben. Die Friedensthematik
    ist Aufgabe nicht nur des AStA, sondern der gesamten Gesellschaft.
    Geoffrey möchte als erstes die Denkweise ändern, wie Referate übergreifend
    zusammenarbeiten können. Er stellt die Frage, wie man wieder eine Universität
    hinbekommt, an der man zusammenarbeitet und nicht in Scheuklappen denkt.

    Karima antwortet für das Referat und möchte die Zusammenarbeit mit den
    teilautonomen Referaten verbessern. Außerdem soll die Arbeit des AStA
    auch außerhalb des AStA stattfinden und in die Universität getragen werden.
    Mario sagt, dass Frieden übergreifend ist und in der Kritischen Wissenschaft
    die Zivilklausel verortet wäre.

    Eine Studentin namens Lena fragt, inwiefern der AStA das Thema
    Feminismus vertritt.

    Golnar fragt, wie sich der AStA zur Frage der Inklusion verhält. Golnar fragt,
    wie alle das Verhältnis von StuPa zu AStA sehen. Sie fragt, welche Debatten
    im Bereich der Kultur anstehen. Golnar fragt alle Referate, wie sich diese
    an der Anti-ExzellenzIni-Kampagne beteiligen möchten.

    Karima beantwortet die Frage zur Diskriminierung. Es kann Personen geben,
    die sowohl diskriminiert werden als auch selber diskriminieren.

    Mario erklärt das Verhältnis von StuPa zu AStA damit, dass die Referent*innen
    vom StuPa bestätigt werden.

    Seld geht auf Frage zu Feminismus ein. Dabei erwähnt er mehrere
    Frauenbewegungen.

    Artur sagt, dass die formale Frage, wer Referent oder Referentin ist,
    nicht relevant ist, da Frauke Petri und Beatrix von Storch beides Frauen sind,
    aber sehr frauenfeindlich agieren.

    Armin und Seld stellen klar, dass es keine Referatsmeinung zu dem
    Thema gibt und es unterschiedliche Positionen im Referat dazu gibt.

    Janis beantwortet die Frage zu den anstehenden Debatten. Beispielsweise
    stellt sich die Frage, ob "tanzen gegen rechts" genug ist.

    Mena geht auf die ExzellenzIni ein.

    23:30

    Beispielsweise geht sie darauf ein, dass ein Studi ihr gesagt hat, dass
    er über die ExIni über das Klo herausgefunden hat.

    Die Bestätigung findet offen statt. 7 der Referent*Innen sind stimmberechtigt
    im StuPa.

    16+7:6:0, Referent*innen damit bestätigt.

    \subsection{Antrag Ältestenrat zur Beauftragung des SWOGA}

    Geoffrey erklärt den Antrag vom Ältestenrat.

    Melf sagt, dass sich der SWOGA am 11. Mai um 18:30 das nächste Mal
    trifft.

    M:0:1 für den Antrag.

    \subsection{Verschiedenes}

    Till erinnert, dass StuPa beim letzten Mal einen Aufruf zur Teilnahme
    an der Erinnerung an das Bücherverbrennen beschlossen hat. Das findet
    statt am Kaiser-Friedrich-Ufer am 10. Mai. Alle sind dazu aufgerufen
    dort hinzugehen.

    Vincent erinnert daran, dass der AStA zusammen mit anderen Gruppen
    für den 1. Mai aufruft. Am 3. Mai findet die Demonstration gegen
    die Schließung der Holzwirtschaft statt. Geht um 15:30 los am Dammtorbahnhof.

    Artur verweist auf eine Veranstaltung von Torsten Bultmann. 10. Mai um 18 Uhr.

    Sitzung um 23:41 geschlossen.

\end{document}
